% =============================================================================
%  Section 6 — Conclusion
% =============================================================================
\section{Conclusion}
\label{sec:conc}

% -----------------------------------------------------------------------------
\subsection{Summary}
\label{sec:conc-summary}

We have argued that multi-tenant LLM agent runtimes introduce a class of
\emph{coupled resource--security failures} that layer-separated isolation
cannot prevent: under resource contention, the security module is itself
starved of the resources it needs to enforce, and during the contention
window that precedes a hard failure, privileged operations slip through a
permission layer that has no causal link to the ongoing attack. The root
cause is the absence of a shared causal identifier across the governance
layers, compounded by a LIM-specific attack surface---spawn-time privilege
escalation, tenant-identity loss across spawn, and depth-oblivious
approval---that classical isolation models were never designed to handle.

To close this gap, we presented \textbf{Neuron}, a multi-tenant LLM agent
governance runtime built on the principle of \emph{coupled governance},
and made four contributions. \emph{First}, we presented a coupled
resource--security governance architecture that fuses cgroup-based
resource control, sandbox-based execution isolation, and permission-based
access control through two complementary mechanisms: a unified audit
chain keyed by a per-turn trace identifier for cross-layer audit
correlation, and a cross-layer joint-decision mechanism in which a
resource-layer signal tightens the permission verdict in real time,
together with a pseudocode-level invariant argument showing that
spawn-time privilege non-increase holds across all code paths.
\emph{Second}, we closed the LIM attack surface through spawn-time
privilege non-increase, tenant-identity propagation, and a depth-aware
escalation policy using a dual-track destructive-tool classification
(name whitelist + capability-surface classifier) that closes the
name-evasion bypass. \emph{Third}, we stated an explicit threat model
and enumerated the conditions under which the guarantees degrade or
fail. \emph{Fourth}, we validated the architecture through reproducible
experiments yielding statistically aggregated evidence across two
core governance scenarios, four adaptive probes, and an
external-framework architectural comparison contrasting Neuron's
coupled governance with the layer-separated designs of AutoGen,
LangGraph, and MCP---noting explicitly that these are general-purpose
orchestration frameworks rather than multi-tenant security systems,
so the comparison characterizes the architectural gap rather than
claims victory over a straw man. The
benign-workload evaluation confirms a $0.00\%$ false-positive rate across
14{,}850 operations and governance overhead below $0.01\%$ of
end-to-end latency.

The evaluation substantiates the coupled governance thesis along three
evidentiary tiers. \emph{First}, a dedicated permission-layer starvation
experiment (Section~\ref{sec:eval-starvation}) directly induces the
central coupled-failure mode: under resource pressure, the permission
decision correctness drops to $66.67\%$ (Baseline) and $63.33\%$
(Permission-only) due to deadline fallback to the permissive default,
while coupled governance preserves $100\%$ correctness with a $6.8\times$
p95 latency improvement---direct evidence that the permission layer
degrades under resource pressure and that coupling prevents the
degradation. \emph{Second}, a cross-layer joint-decision experiment
(Section~\ref{sec:eval-joint}) shows a resource-layer signal
(rate-limit rejection count) \emph{changing} the permission verdict in
real time: under sustained pressure, the depth threshold tightens from
$\geq 2$ to $\geq 1$, intercepting depth~1 destructive-tool requests
that the static policy would have allowed---the missing ``resource
signal actually changed the permission decision'' scenario that
distinguishes coupled governance from three independent modules sharing
a trace identifier. \emph{Third}, the original Scenario~6 results hold:
under a simultaneous resource--permission cross-attack, coupled
governance achieves a $100\%$ interception rate across all three attack
vectors, where permission-only enforcement achieves only $\approx 67\%$
and the baseline $0\%$; under four adaptive attacks, it remains
effective against trace forgery, ThreadLocal leakage, and
depth-spoofing, with the residual low-and-slow surface disclosed in
the failure-case analysis. The original $0\%$ attack success rate
was conditioned on a
destructive-tool name whitelist; integrating the capability-surface
classifier (Section~\ref{sec:eval-capability-classifier}) closes the
F6 name-evasion bypass, raising the interception rate on
name-evasion attacks from $0\%$ (name-only) to $100\%$
(capability-integrated) with a $0\%$ false-positive rate.

We emphasize that the Scenario~6 policy is a deterministic pure function
of the seeded inputs, so its $0\%$ attack success rate is reported as a
raw proportion without inferential intervals; the genuinely stochastic
binary metrics are the supplementary real-LLM experiments
(Section~\ref{sec:eval-supplement}) and the external-framework
real-AutoGen reproduction (Section~\ref{sec:eval-external}), where the
Wilson CI serves its standard inferential role. Under
resource-contention flooding, coupled governance confines co-resident
normal-agent read latency to the uncontended regime
(p95\,$\approx$\,0.029\,ms), a $11.8\times$ improvement at p95 and
$20.4\times$ at p99 over the layer-separated baseline, with all
latencies directly measured via \texttt{System.nanoTime()}. The same
governance gap is reproduced on AutoGen, LangGraph, and MCP. These are
general-purpose orchestration frameworks rather than multi-tenant
security systems; no directly comparable baseline exists. The
architectural differentiation from prior systems is threefold: relative
to AIOS~\cite{mei2025aios}, Neuron couples rather than isolates the
resource and permission layers; relative to the dual-LLM
pattern~\cite{anthropic2024dual}, Neuron enforces privilege
non-increase at the runtime layer rather than relying on a trusted
planner LLM; relative to classical container isolation, Neuron treats
the LLM token and spawn tree as first-class governance objects. The
coupled-failure surface is the contribution. This is attained at a
measured $0.00\%$ false-positive rate across 14{,}850 benign
operations on a live multi-tenant workload, and the governance
overhead is a small fraction of a single LLM call
(Section~\ref{sec:eval-e2e}). The coupling of resource and security
enforcement---not either layer in isolation---is what closes the
contention-window attack surface, establishing coupled governance as
an architectural baseline for multi-tenant agent runtimes.

% -----------------------------------------------------------------------------
\subsection{Limitations}
\label{sec:conc-lim}

The failure-case analysis of Section~\ref{sec:lim-failure} identifies
eight conditions under which the coupled governance guarantees degrade
or fail. We restate the most consequential of them here, in condensed
form, as the explicit limitations of this work.

\paragraph{L1: Low-and-slow attacks under the rate threshold.}
The rate limiter is rate-based, not morphology-based. An attacker operating
at 19 writes/s---just below the 20 writes/s per-tenant threshold---bypasses
the limiter by construction (Section~\ref{sec:eval-adaptive}, A1). Benign
impact is bounded (p95\,$<$\,0.1\,ms) but not eliminated. This is the
fundamental limitation of rate-based governance.

\paragraph{L2: Prevention-layer scope.}
The threat model of Section~\ref{sec:threat} assumes the adversary can
already issue malicious escalation requests; Neuron's defense is structural
enforcement at the runtime layer, not prevention at the prompt-injection
layer. Indirect prompt injection that succeeds in coercing the agent to
\emph{issue} a request is out of scope; the runtime's contribution is to
ensure that any such request, once issued, is adjudicated under
spawn-tree context and resource pressure.

\paragraph{L3: Single-node scope.}
The core evaluation uses a single runtime instance. A supplementary
multi-node experiment (Section~\ref{sec:eval-distributed}) verifies that
governance properties---trace propagation, privilege non-increase,
depth-aware escalation, joint-decision coupling, and audit
correlation---hold across three nodes via HTTP federation, with
$30/30$ correctness on all five scenarios and cross-node overhead of
$2$--$3\times$ (single-digit milliseconds). However, distributed
resource-pressure consensus (which would require the Raft path) and
cross-node tenant-identity attestation via cryptographic credentials
remain open; the current federated experiment propagates governance
state but does not achieve consensus on shared resource signals.

\paragraph{L4: Sub-millisecond improvement under LLM-dominant regimes.}
The $11.8\times$ p95 improvement is measured at the VFS layer, where
latencies are sub-millisecond. In end-to-end turns dominated by LLM
inference (200\,ms--2\,s per call), the governance overhead is a small
fraction of total latency (Section~\ref{sec:eval-e2e}). The improvement
remains significant for high-throughput multi-tenant runtimes---where
contended VFS access is the bottleneck---but is not the dominant factor
in single-turn, LLM-bound workloads.

\paragraph{L5: Partial formal verification of the spawn invariant.}
The invariant argument of Section~\ref{sec:arch-invariant} is
pseudocode and enumeration-based. We supplemented it with a TLA$^+$
specification and TLC model checking
(Section~\ref{sec:eval-tla-scale}): the three spawn invariants
(privilege non-increase, no-leak-on-exception, no-cross-tenant-leak)
are verified over 1{,}681 reachable states in a finite model
(3~privilege levels $\times$ 2~tenants $\times$ 2~spawn depths $\times$
2~concurrent threads), and three mutation configurations each
precisely violate the target invariant, confirming the specification
is non-vacuous. This is a partial verification---not a
machine-checked proof in the style of seL4~\cite{klein2009sel4}---and
a future full mechanization in a proof assistant (Coq or Lean) is
needed to close the gap.

\paragraph{L6: Evaluation scale.}
The security and false-positive evaluation uses 3~tenants $\times$
3~agents $=$ 9~concurrent agents over 3--10~second windows. A
dedicated scalability benchmark (Section~\ref{sec:eval-tla-scale}) extends
this to 150~tenants $\times$ 3~agents $=$ 450~concurrent agents and
confirms that the two structural governance primitives---audit-chain
append and \texttt{SpawnPrivilegeContext} propagation---remain stable
across the $50\times$ concurrency range (616--1545\,ops/s and
1.6--3.1\,$\mu$s respectively), and that per-tenant token-bucket
partitioning prevents lock contention from growing with $N$. At
$N \geq 100$ (300+ agents), the writeText measurement saturates
because the rate limiter denies most writes, but the governance
primitives continue to function. The denial-path audit I/O remains an
engineering optimization target (batch appends, asynchronous denial
logging) rather than a structural flaw. The $50\times$ scaling data
substantially narrows the uncertainty for $\geq 100$~tenant deployment.

\paragraph{L7: Heuristic coupling rule.}
The cross-layer joint-decision rule of Section~\ref{sec:eval-joint}---a
hand-tuned threshold (\texttt{PRESSURE\_THRESHOLD}=50) above which the
depth threshold tightens by one level---is a heuristic calibrated on
the depth-based escalation morphology evaluated in this paper. It is
not a learned policy, nor is it provably optimal: a different attack
morphology (e.g.\ breadth-based fan-out, or a non-depth escalation
axis) would require a different coupling rule and threshold
recalibration. The threshold selection is justified empirically by the
token-bucket geometry and the deadline-fallback latency
(Section~\ref{sec:eval-joint}), but the rule's generalization to
unseen morphologies is unproven. A learned coupling function that
adapts the tightening rule to the observed attack morphology is left
as future work.

\paragraph{L8: Single-system evaluation (partially addressed).}
All governance-effectiveness, latency, and scalability measurements
are conducted on Neuron's own Java/JVM implementation. The
external-framework comparison (Section~\ref{sec:eval-external})
validates that the attack surface is structural---the same attacks
succeed on AutoGen, LangGraph, and MCP---and the cross-system porting
experiment (Section~\ref{sec:eval-external}, ``Cross-system porting
of the joint-decision mechanism'') further validates that the core
coupling mechanism (\texttt{ResourcePressureAwareEscalationPolicy})
can be ported to real AutoGen~0.7.5 as a framework-agnostic
middleware, producing the same resource-signal-changes-permission-verdict
behavior ($100\%$ interception on AutoGen + joint decision vs.\ $0\%$
on AutoGen + static). However, the portability of the full coupling
mechanism (trace identifier injection, depth-aware escalation, and
pressure-aware joint decision together) to non-JVM runtimes
(Python, Rust, Go) and to non-cgroup resource backends (container
orchestration platforms, serverless schedulers) is not yet tested.
We disclose the residual gap: the contribution is the governance
\emph{design} validated on one full realization (Neuron) plus one
ported mechanism (joint decision on AutoGen), not a portable library
demonstrated across multiple complete runtimes.

% -----------------------------------------------------------------------------
\subsection{Future Work}
\label{sec:conc-future}

Three extensions follow naturally from the limitations above.

First, we will introduce \emph{class-aware \IO{} rate limiting} that
discriminates not only between reads and writes (as the current
\texttt{VfsRateLimiter} does) but also between \IO{} \emph{classes}---metadata
lookups, bulk scans, embedding-store fetches, and network egress---each with
its own token bucket and its own coupling to the permission layer. Class-aware
limiting would let the runtime throttle the specific \IO{} morphology of a
low-and-slow attack (L1) while leaving benign mixed workloads unaffected,
sharpening the source-side throttling on which the coupled enforcement loop
depends.

Second, we will extend coupled governance from a single runtime to a
\emph{multi-node cluster} setting, in which the per-turn trace identifier
spans runtime boundaries and the unified audit chain is aggregated across
nodes. Distributed coupled governance raises new questions---cross-node
tenant-identity attestation, consensus on resource-pressure signals, and a
federated escalation policy---that the single-node \texttt{TraceContext}
does not address (L3). It is the prerequisite for governing multi-tenant
agent fleets that share GPUs and model-serving front-ends across machines.

Third, we will \emph{mechanize the spawn-time privilege non-increase
invariant} in a proof assistant (e.g., Coq or Lean), targeting a
machine-checked proof that the property holds for all reachable states of
the spawn path (L5). The pseudocode-level argument of
Section~\ref{sec:arch-invariant} provides the proof outline; mechanization
will require formalizing the \texttt{PermissionProfile} subsumption order,
the \texttt{SpawnPrivilegeContext} semantics, and the virtual-thread
inheritance contract, and discharging the concurrent-case obligations that
are currently argued informally.
